\documentclass[12pt,a4paper,twocolumn]{article}

% ==================== CONFIGURACIÓN DE IDIOMAS ====================
\usepackage[utf8]{inputenc}
\usepackage[T1]{fontenc}
\usepackage[spanish]{babel}
\usepackage{csquotes}

% ==================== PAQUETES MATEMÁTICOS ====================
\usepackage{amsmath,amssymb,amsthm,mathtools}
\usepackage{physics} % Para derivadas y notación física
\usepackage{siunitx} % Para unidades SI

% ==================== PAQUETES DE GRÁFICOS ====================
\usepackage{graphicx}
\usepackage{subcaption}
\usepackage{float}
\usepackage{tikz}
\usepackage{pgfplots}
\pgfplotsset{compat=1.18}

% ==================== PAQUETES DE CÓDIGO ====================
\usepackage{listings}
\usepackage{xcolor}

% Configuración de listings para Python
\lstset{
    language=Python,
    basicstyle=\ttfamily\small,
    keywordstyle=\color{blue}\bfseries,
    commentstyle=\color{green!60!black},
    stringstyle=\color{red},
    numbers=left,
    numberstyle=\tiny\color{gray},
    stepnumber=1,
    numbersep=5pt,
    backgroundcolor=\color{gray!10},
    frame=single,
    breaklines=true,
    breakatwhitespace=true,
    tabsize=4,
    showstringspaces=false
}

% ==================== PAQUETES DE TABLAS ====================
\usepackage{booktabs}
\usepackage{multirow}
\usepackage{array}
\usepackage{longtable}

% ==================== PAQUETES DE REFERENCIAS ====================
\usepackage{hyperref}
\hypersetup{
    colorlinks=true,
    linkcolor=blue,
    filecolor=magenta,      
    urlcolor=cyan,
    citecolor=red,
    pdftitle={Sistema de Análisis Acústico de Flautas Traverso},
    pdfauthor={Autor},
    pdfsubject={Análisis Acústico de Instrumentos Musicales},
    pdfkeywords={flauta traverso, OpenWind, impedancia acústica, análisis numérico}
}

% ==================== CONFIGURACIÓN ADICIONAL ====================
\usepackage{geometry}
\geometry{margin=2.5cm}
\usepackage{fancyhdr}
\usepackage{titlesec}
\usepackage{enumitem}
\usepackage{multicol}

% ==================== COMANDOS PERSONALIZADOS ====================
\newcommand{\openwind}{\textsc{OpenWind}}
\newcommand{\flutees}{flauta traverso}
\newcommand{\impedance}{Z}
\newcommand{\admittance}{Y}
\newcommand{\freq}{f}
\newcommand{\freqzero}{f_0}
\newcommand{\freqone}{f_1}
\newcommand{\freqplay}{f_{\text{play}}}
\newcommand{\freqtwo}{f_2}

% ==================== METADATOS ====================
\title{\textbf{Sistema de Análisis Acústico y Visualización\\de Flautas Traverso}\\
\large{Integración de OpenWind para Análisis Numérico\\y Caracterización Acústica}}
\author{Autor del Sistema}
\date{\today}

% ==================== INICIO DEL DOCUMENTO ====================
\begin{document}

\maketitle

% ==================== RESUMEN EJECUTIVO ====================
\begin{abstract}
Este documento presenta un sistema completo de análisis acústico y visualización para flautas traverso históricas, desarrollado mediante la integración de la biblioteca \openwind{} con herramientas de visualización 2D/3D y gestión de datos. El sistema permite la caracterización acústica sistemática mediante el cálculo de impedancia, análisis de inharmonicidad, compresión modal de octava (MOC), y múltiples métricas adicionales. Se incluyen derivaciones matemáticas completas de todas las mediciones, explicación detallada de la integración con \openwind{}, y casos de estudio con instrumentos históricos.

\vspace{0.3cm}
\textbf{Palabras clave:} flauta traverso, \openwind{}, impedancia acústica, análisis numérico, caracterización instrumental
\end{abstract}

\tableofcontents
\newpage

% ==================== SECCIÓN 1: INTRODUCCIÓN ====================
\section{Introducción}

\subsection{Contexto Histórico y Musical}

Las flautas traverso constituyen una familia de instrumentos de viento madera que han evolucionado significativamente desde el Renacimiento hasta la actualidad. Durante el período barroco (siglos XVII-XVIII), estos instrumentos alcanzaron un desarrollo técnico y artístico notable, con constructores como Hotteterre, Quantz y otros estableciendo estándares de diseño que influenciaron generaciones posteriores.

La caracterización acústica sistemática de estos instrumentos históricos presenta desafíos únicos:
\begin{itemize}
    \item Variabilidad en materiales y técnicas de construcción
    \item Necesidad de análisis no destructivo
    \item Requerimiento de comparación objetiva entre instrumentos
    \item Relación compleja entre geometría y comportamiento acústico
\end{itemize}

\subsection{Objetivos del Sistema}

El sistema desarrollado tiene como objetivos principales:

\begin{enumerate}
    \item \textbf{Análisis Acústico Sistemático}: Cálculo de impedancia acústica mediante métodos numéricos validados
    \item \textbf{Caracterización Completa}: Múltiples métricas acústicas para evaluación objetiva
    \item \textbf{Visualización Integral}: Representación 2D y 3D de geometría y resultados acústicos
    \item \textbf{Gestión de Datos}: Base de datos para almacenamiento y recuperación eficiente
    \item \textbf{Herramientas de Diseño}: Editor de geometría y generación de planos de ingeniería
\end{enumerate}

% ==================== SECCIÓN 2: ARQUITECTURA DEL SISTEMA ====================
\section{Arquitectura del Sistema}

\subsection{Visión General}

El sistema está estructurado en módulos independientes pero interconectados, siguiendo principios de diseño orientado a objetos y separación de responsabilidades. La Figura~\ref{fig:arquitectura} muestra el diagrama de arquitectura general.

\begin{figure}[H]
\centering
% Aquí iría el diagrama de arquitectura (usar TikZ o incluir imagen)
\caption{Arquitectura general del sistema}
\label{fig:arquitectura}
\textit{[Incluir diagrama de arquitectura aquí]}
\end{figure}

\subsection{Componentes Principales}

\subsubsection{FluteData y FluteDataDB}

La clase \texttt{FluteData} constituye el núcleo del sistema, responsable de:
\begin{itemize}
    \item Carga y validación de datos geométricos desde archivos JSON
    \item Concatenación de mediciones de múltiples partes (headjoint, left, right, foot)
    \item Normalización de coordenadas (corcho en $x=0$)
    \item Preparación de datos para \openwind{}
\end{itemize}

\texttt{FluteDataDB} extiende \texttt{FluteData} agregando:
\begin{itemize}
    \item Integración con base de datos SQLite
    \item Caché de resultados de cálculo
    \item Detección automática de cambios en parámetros
    \item Recuperación eficiente de resultados previos
\end{itemize}

\subsubsection{FluteOperations}

Proporciona métodos estáticos para:
\begin{itemize}
    \item Visualización 2D de perfiles (físico y acústico)
    \item Generación de gráficos de análisis acústico
    \item Cálculo de métricas derivadas
    \item Operaciones de geometría (concatenación, interpolación)
\end{itemize}

\subsubsection{FluteAnalyzer}

Encapsula todos los análisis acústicos:
\begin{itemize}
    \item Cálculo de inharmonicidad
    \item Cálculo de MOC
    \item Cálculo de B$_I$ y ESPE
    \item Generación de gráficos comparativos
\end{itemize}

\subsubsection{Interfaz Gráfica Unificada}

Desarrollada en PyQt5, proporciona:
\begin{itemize}
    \item Carga y selección de flautas
    \item Visualización interactiva 2D/3D
    \item Análisis acústico con múltiples métricas
    \item Generación de planos de ingeniería
    \item Editor de geometría interactivo
    \item Generación de código G para CNC
\end{itemize}

% ==================== SECCIÓN 3: INTEGRACIÓN CON OPENWIND ====================
\section{Integración con OpenWind}

\subsection{Introducción a OpenWind}

\openwind{} es una biblioteca Python desarrollada por el equipo de acústica del INRIA (Institut National de Recherche en Informatique et en Automatique) para el modelado numérico de instrumentos de viento. Implementa métodos de elementos finitos y diferencias finitas para resolver las ecuaciones de la acústica lineal en conductos de sección variable.

\subsection{Modelo Físico}

\openwind{} resuelve las ecuaciones de la acústica lineal en el dominio de la frecuencia. Para un conducto de sección variable $S(x)$, las ecuaciones fundamentales son:

\subsubsection{Ecuaciones de Continuidad y Momentum}

La ecuación de continuidad para ondas acústicas en un conducto es:
\begin{equation}
\frac{\partial (\rho' S)}{\partial t} + \frac{\partial (\rho_0 S u)}{\partial x} = 0
\label{eq:continuidad}
\end{equation}

donde $\rho'$ es la densidad acústica, $\rho_0$ la densidad de equilibrio, $S(x)$ la sección transversal, y $u$ la velocidad acústica.

La ecuación de momentum (conservación de cantidad de movimiento) es:
\begin{equation}
\rho_0 \frac{\partial u}{\partial t} + \frac{\partial p}{\partial x} = 0
\label{eq:momentum}
\end{equation}

donde $p$ es la presión acústica.

\subsubsection{Ecuación de Estado}

Para un gas ideal con procesos adiabáticos:
\begin{equation}
p = c_0^2 \rho'
\label{eq:estado}
\end{equation}

donde $c_0$ es la velocidad del sonido:
\begin{equation}
c_0 = \sqrt{\gamma \frac{P_0}{\rho_0}}
\label{eq:velocidad_sonido}
\end{equation}

con $\gamma$ el coeficiente adiabático ($\gamma \approx 1.4$ para aire) y $P_0$ la presión de equilibrio.

\subsubsection{Ecuación de Helmholtz}

Combinando las ecuaciones anteriores y asumiendo dependencia temporal armónica $e^{j\omega t}$, se obtiene la ecuación de Helmholtz para el conducto:

\begin{equation}
\frac{d^2 p}{dx^2} + \frac{1}{S}\frac{dS}{dx}\frac{dp}{dx} + k^2 p = 0
\label{eq:helmholtz}
\end{equation}

donde $k = \omega/c_0$ es el número de onda y $\omega = 2\pi f$ es la frecuencia angular.

\subsection{Configuración del Modelo para Flauta Traversa}

\subsubsection{Player("FLUTE")}

La clase \texttt{Player} en \openwind{} configura parámetros específicos del instrumento. Para flauta traversa, \texttt{Player("FLUTE")} establece:

\begin{itemize}
    \item \textbf{Radiación de embocadura}: Modelo de radiación desde el agujero de embocadura hacia el exterior
    \item \textbf{Radiación de agujeros}: Modelo \textit{unflanged} para agujeros laterales abiertos
    \item \textbf{Radiación de campana}: Modelo \textit{unflanged} para el extremo abierto
    \item \textbf{Pérdidas en paredes}: Modelo de Webster-Lokshin para pérdidas viscales y térmicas
\end{itemize}

\subsubsection{Construcción de la Geometría}

La geometría se construye a partir de mediciones discretas del diámetro interno. Para cada parte de la flauta, se tienen puntos $(x_i, d_i)$ donde $x_i$ es la posición axial y $d_i$ el diámetro interno.

La conversión a formato \openwind{} requiere:
\begin{enumerate}
    \item \textbf{Normalización de coordenadas}: El corcho (stopper) se coloca en $x=0$
    \item \textbf{Conversión a radios}: $r_i = d_i/2$
    \item \textbf{Conversión de unidades}: De milímetros a metros
    \item \textbf{Interpolación}: \openwind{} interpola entre puntos discretos
\end{enumerate}

El formato de geometría para \openwind{} es:
\begin{equation}
\text{geom} = [[x_0, r_0], [x_1, r_1], \ldots, [x_n, r_n]]
\label{eq:geometria_openwind}
\end{equation}

donde todas las coordenadas están en metros.

\subsubsection{Agujeros Laterales}

Los agujeros laterales se especifican mediante:
\begin{itemize}
    \item Posición axial $x_h$
    \item Diámetro $d_h$
    \item Altura de chimenea $h_h$ (chimney height)
    \item Estado (abierto/cerrado según digitación)
\end{itemize}

Para cada nota, se especifica qué agujeros están abiertos mediante un archivo de digitación (fingering chart).

\subsubsection{Cálculo de Impedancia}

El objeto \texttt{ImpedanceComputation} calcula la impedancia acústica de entrada:

\begin{equation}
\impedance(f) = \frac{p(0)}{U(0)}
\label{eq:impedancia}
\end{equation}

donde $p(0)$ es la presión acústica en la embocadura y $U(0) = S(0) u(0)$ es el flujo volumétrico acústico.

La admitancia es:
\begin{equation}
\admittance(f) = \frac{1}{\impedance(f)}
\label{eq:admittancia}
\end{equation}

\openwind{} resuelve la ecuación de Helmholtz numéricamente para un rango de frecuencias:
\begin{equation}
f \in [f_{\min}, f_{\max}] \quad \text{con paso } \Delta f
\label{eq:rango_frecuencias}
\end{equation}

En nuestro sistema: $f_{\min} = \SI{100}{\hertz}$, $f_{\max} = \SI{3000}{\hertz}$, $\Delta f = \SI{2}{\hertz}$.

\subsubsection{Frecuencias de Antiresonancia}

Las frecuencias de antiresonancia corresponden a los máximos de admitancia (mínimos de impedancia). Se extraen mediante:

\begin{equation}
\freqzero = \arg\max_f |\admittance(f)|
\label{eq:antiresonancia_1}
\end{equation}

\begin{equation}
\freqone = \arg\max_{f > \freqzero} |\admittance(f)|
\label{eq:antiresonancia_2}
\end{equation}

Estas frecuencias son fundamentales para el cálculo de todas las métricas acústicas.

% ==================== SECCIÓN 4: MEDICIONES ACÚSTICAS ====================
\section{Mediciones Acústicas}

\subsection{Inharmonicidad}

\subsubsection{Definición y Derivación}

La inharmonicidad mide la desviación del segundo pico de impedancia respecto al doble del primero. Para un sistema perfectamente armónico, esperaríamos:

\begin{equation}
\freqone = 2\freqzero
\label{eq:armonico_ideal}
\end{equation}

Sin embargo, en instrumentos reales, esta relación no se cumple exactamente. La inharmonicidad se define como:

\begin{equation}
\text{Inharmonicidad} = 1200 \log_2\left(\frac{\freqone}{2\freqzero}\right) \quad \text{[cents]}
\label{eq:inharmonicidad}
\end{equation}

\subsubsection{Derivación Completa}

Partiendo de la definición de cents:
\begin{equation}
\text{cents} = 1200 \log_2\left(\frac{f_2}{f_1}\right)
\label{eq:cents_def}
\end{equation}

Para nuestro caso, comparamos $\freqone$ con $2\freqzero$:
\begin{align}
\text{Inharmonicidad} &= 1200 \log_2\left(\frac{\freqone}{2\freqzero}\right) \\
&= 1200 \left[\log_2(\freqone) - \log_2(2\freqzero)\right] \\
&= 1200 \left[\log_2(\freqone) - \log_2(2) - \log_2(\freqzero)\right] \\
&= 1200 \left[\log_2(\freqone) - 1 - \log_2(\freqzero)\right] \\
&= 1200 \log_2\left(\frac{\freqone}{\freqzero}\right) - 1200
\label{eq:inharmonicidad_deriv}
\end{align}

\subsubsection{Interpretación Física}

\begin{itemize}
    \item \textbf{Inharmonicidad = 0 cents}: Sistema perfectamente armónico
    \item \textbf{Inharmonicidad > 0}: Segundo pico más alto de lo esperado (sobretonos más agudos)
    \item \textbf{Inharmonicidad < 0}: Segundo pico más bajo de lo esperado (sobretonos más graves)
\end{itemize}

La inharmonicidad está relacionada con:
\begin{itemize}
    \item Efectos de extremos abiertos
    \item Comportamiento no lineal del flujo
    \item Acoplamiento entre modos
\end{itemize}

\subsection{MOC (Modal Octave Compression)}

\subsubsection{Definición y Contexto Físico}

La Compresión Modal de Octava (MOC) es una métrica que cuantifica el comportamiento no lineal del instrumento entre la primera y segunda octava. Fue introducida por Benade y otros investigadores para caracterizar cómo se comportan los modos acústicos en instrumentos de viento.

\subsubsection{Derivación Completa}

Partimos de la relación entre frecuencia y longitud efectiva para un tubo abierto en ambos extremos:

\begin{equation}
f_n = \frac{n c_0}{2 L_{\text{eff}}}
\label{eq:frecuencia_tubo}
\end{equation}

donde $n$ es el número del modo, $c_0$ la velocidad del sonido, y $L_{\text{eff}}$ la longitud efectiva del tubo.

Para la primera octava ($n=1$):
\begin{equation}
\freqplay = \frac{c_0}{2 L_{\text{eff},I}}
\label{eq:freq_play_I}
\end{equation}

Para la segunda octava ($n=2$):
\begin{equation}
2\freqplay = \frac{c_0}{L_{\text{eff},II}}
\label{eq:freq_play_II}
\end{equation}

Sin embargo, las frecuencias medidas $\freqzero$ y $\freqone$ no coinciden exactamente con estas frecuencias teóricas debido a efectos de extremos y agujeros.

Definimos las desviaciones:
\begin{align}
\Delta L_I &= L_{\text{eff},I} - L_{\text{real}} = \frac{c_0}{2}\left(\frac{1}{\freqplay} - \frac{1}{\freqzero}\right) \\
\Delta L_{II} &= L_{\text{eff},II} - L_{\text{real}} = c_0\left(\frac{1}{2\freqplay} - \frac{1}{\freqone}\right)
\label{eq:desviaciones_longitud}
\end{align}

La diferencia entre estas desviaciones es:
\begin{equation}
\Delta\Delta L = \Delta L_{II} - \Delta L_I = c_0\left[\left(\frac{1}{2\freqplay} - \frac{1}{\freqone}\right) - \left(\frac{1}{\freqplay} - \frac{1}{\freqzero}\right)\right]
\label{eq:delta_delta_L}
\end{equation}

Simplificando:
\begin{align}
\Delta\Delta L &= c_0\left[\frac{1}{2\freqplay} - \frac{1}{\freqone} - \frac{1}{\freqplay} + \frac{1}{\freqzero}\right] \\
&= c_0\left[\frac{1}{\freqzero} - \frac{1}{\freqone} - \frac{1}{2\freqplay}\right] \\
&= c_0\left[\left(\frac{1}{\freqzero} - \frac{1}{\freqplay}\right) - \left(\frac{1}{\freqone} - \frac{1}{2\freqplay}\right)\right]
\label{eq:delta_delta_L_simplificado}
\end{align}

El MOC se define como la razón entre estas diferencias de períodos:
\begin{equation}
\text{MOC} = \frac{\frac{1}{\freqone} - \frac{1}{2\freqplay}}{\frac{1}{\freqzero} - \frac{1}{\freqplay}}
\label{eq:moc_def}
\end{equation}

\subsubsection{Interpretación Física}

\begin{itemize}
    \item \textbf{MOC = 1.0}: Compresión ideal, comportamiento lineal perfecto
    \item \textbf{MOC > 1.0}: Sobre-compresión, la segunda octava está más comprimida
    \item \textbf{MOC < 1.0}: Sub-compresión, la segunda octava está menos comprimida
\end{itemize}

El MOC está relacionado con:
\begin{itemize}
    \item Efectos de extremos abiertos
    \item Comportamiento de agujeros laterales
    \item Acoplamiento entre modos acústicos
    \item Efectos no lineales del flujo
\end{itemize}

\subsection{B$_I$ (Bias de la Primera Octava)}

\subsubsection{Definición}

B$_I$ mide la desviación de la frecuencia de resonancia medida $\freqzero$ respecto a la frecuencia teórica de ejecución $\freqplay$:

\begin{equation}
B_I = 1200 \log_2\left(\frac{\freqplay}{\freqzero}\right) \quad \text{[cents]}
\label{eq:bi_def}
\end{equation}

\subsubsection{Derivación desde la Teoría de Longitud Efectiva}

Para un tubo abierto en ambos extremos, la frecuencia fundamental es:
\begin{equation}
\freqplay = \frac{c_0}{2 L_{\text{eff}}}
\label{eq:freq_fundamental}
\end{equation}

donde $L_{\text{eff}}$ incluye correcciones de extremos:
\begin{equation}
L_{\text{eff}} = L_{\text{geom}} + \Delta L
\label{eq:longitud_efectiva}
\end{equation}

La corrección $\Delta L$ se debe a:
\begin{itemize}
    \item Efectos de extremo abierto (end correction)
    \item Efectos de agujeros abiertos
    \item Acoplamiento con el exterior
\end{itemize}

Si la frecuencia medida $\freqzero$ difiere de $\freqplay$, esto indica una longitud efectiva diferente:
\begin{equation}
\freqzero = \frac{c_0}{2 L_{\text{eff},0}}
\label{eq:freq_medida}
\end{equation}

La relación entre frecuencias es:
\begin{equation}
\frac{\freqplay}{\freqzero} = \frac{L_{\text{eff},0}}{L_{\text{eff}}}
\label{eq:relacion_frecuencias}
\end{equation}

Tomando logaritmo en base 2 y multiplicando por 1200:
\begin{align}
B_I &= 1200 \log_2\left(\frac{\freqplay}{\freqzero}\right) \\
&= 1200 \log_2\left(\frac{L_{\text{eff},0}}{L_{\text{eff}}}\right)
\label{eq:bi_deriv}
\end{align}

\subsubsection{Interpretación}

\begin{itemize}
    \item \textbf{$B_I = 0$ cents}: Afinación perfecta, $\freqzero = \freqplay$
    \item \textbf{$B_I > 0$ cents}: Nota suena más aguda de lo esperado, $\freqzero < \freqplay$
    \item \textbf{$B_I < 0$ cents}: Nota suena más grave de lo esperado, $\freqzero > \freqplay$
\end{itemize}

\subsection{ESPE (Effective Sound Path Extension)}

\subsubsection{Definición}

ESPE cuantifica la diferencia en la extensión efectiva del camino del sonido entre la primera y segunda octava:

\begin{equation}
\text{ESPE} = 1200 \log_2\left(\frac{L_{\text{eff},I}}{L_{\text{eff},I} + \Delta\Delta L}\right) \quad \text{[cents]}
\label{eq:espe_def}
\end{equation}

\subsubsection{Derivación Completa}

De la ecuación~\eqref{eq:freq_play_I}, la longitud efectiva para la primera octava es:
\begin{equation}
L_{\text{eff},I} = \frac{c_0}{2 \freqplay}
\label{eq:longitud_efectiva_I}
\end{equation}

La desviación de longitud para la primera octava es:
\begin{equation}
\Delta L_I = \frac{c_0}{2}\left(\frac{1}{\freqplay} - \frac{1}{\freqzero}\right)
\label{eq:delta_L_I}
\end{equation}

Para la segunda octava:
\begin{equation}
\Delta L_{II} = c_0\left(\frac{1}{2\freqplay} - \frac{1}{\freqone}\right)
\label{eq:delta_L_II}
\end{equation}

La diferencia entre desviaciones es:
\begin{align}
\Delta\Delta L &= \Delta L_{II} - \Delta L_I \\
&= c_0\left[\left(\frac{1}{2\freqplay} - \frac{1}{\freqone}\right) - \left(\frac{1}{\freqplay} - \frac{1}{\freqzero}\right)\right] \\
&= c_0\left[\frac{1}{\freqzero} - \frac{1}{\freqone} - \frac{1}{2\freqplay}\right]
\label{eq:delta_delta_L_espe}
\end{align}

La longitud efectiva corregida para la segunda octava es:
\begin{equation}
L_{\text{eff},II} = L_{\text{eff},I} + \Delta\Delta L
\label{eq:longitud_efectiva_II}
\end{equation}

ESPE mide el cambio relativo:
\begin{align}
\text{ESPE} &= 1200 \log_2\left(\frac{L_{\text{eff},I}}{L_{\text{eff},II}}\right) \\
&= 1200 \log_2\left(\frac{L_{\text{eff},I}}{L_{\text{eff},I} + \Delta\Delta L}\right)
\label{eq:espe_deriv}
\end{align}

\subsubsection{Interpretación Física}

ESPE está relacionado con:
\begin{itemize}
    \item Cambios en la corrección de extremos entre octavas
    \item Comportamiento diferente de agujeros en diferentes frecuencias
    \item Efectos de acoplamiento modal
\end{itemize}

\subsection{Q-Factor (Factor de Calidad)}

\subsubsection{Definición}

El factor de calidad Q mide la agudeza de una resonancia:

\begin{equation}
Q = \frac{\freqzero}{\Delta f}
\label{eq:q_factor}
\end{equation}

donde $\Delta f$ es el ancho de banda a -3 dB (mitad de potencia).

\subsubsection{Derivación desde la Teoría de Resonadores}

Para un resonador con pérdidas, la admitancia cerca de la resonancia puede modelarse como:

\begin{equation}
\admittance(f) \approx \frac{\admittance_{\max}}{1 + j Q \left(\frac{f}{\freqzero} - \frac{\freqzero}{f}\right)}
\label{eq:admittance_resonator}
\end{equation}

El módulo de la admitancia es:
\begin{equation}
|\admittance(f)| = \frac{\admittance_{\max}}{\sqrt{1 + Q^2 \left(\frac{f}{\freqzero} - \frac{\freqzero}{f}\right)^2}}
\label{eq:admittance_modulo}
\end{equation}

A -3 dB (mitad de potencia), tenemos:
\begin{equation}
|\admittance(f)| = \frac{\admittance_{\max}}{\sqrt{2}}
\label{eq:admittance_3db}
\end{equation}

Por lo tanto:
\begin{equation}
1 + Q^2 \left(\frac{f}{\freqzero} - \frac{\freqzero}{f}\right)^2 = 2
\label{eq:q_ecuacion}
\end{equation}

Resolviendo para $f$ cerca de $\freqzero$:
\begin{equation}
Q^2 \left(\frac{f - \freqzero}{\freqzero}\right)^2 \approx 1
\label{eq:q_aproximacion}
\end{equation}

\begin{equation}
\frac{f - \freqzero}{\freqzero} \approx \pm \frac{1}{Q}
\label{eq:q_relacion}
\end{equation}

El ancho de banda es:
\begin{equation}
\Delta f = 2 \frac{\freqzero}{Q}
\label{eq:ancho_banda}
\end{equation}

Por lo tanto:
\begin{equation}
Q = \frac{\freqzero}{\Delta f/2} = \frac{2\freqzero}{\Delta f}
\label{eq:q_final}
\end{equation}

En la práctica, medimos $\Delta f$ como la diferencia entre las frecuencias donde $|\admittance|$ cae a $\admittance_{\max}/\sqrt{2}$.

\subsubsection{Interpretación}

\begin{itemize}
    \item \textbf{Q alto} ($Q > 50$): Resonancia estrecha, color tonal más puro, menos amortiguamiento
    \item \textbf{Q medio} ($10 < Q < 50$): Resonancia moderada, balance entre pureza y riqueza
    \item \textbf{Q bajo} ($Q < 10$): Resonancia ancha, color tonal más rico, más amortiguamiento
\end{itemize}

\subsection{Frecuencias de Resonancia vs Temperamento Igual}

\subsubsection{Definición}

Esta métrica compara las frecuencias de resonancia medidas con las frecuencias del temperamento igual (equal temperament):

\begin{equation}
\text{Desviación} = 1200 \log_2\left(\frac{\freqzero}{\freq_{\text{temp}}}\right) \quad \text{[cents]}
\label{eq:deviation_temperament}
\end{equation}

\subsubsection{Derivación del Temperamento Igual}

En el temperamento igual, la frecuencia de una nota se calcula como:

\begin{equation}
\freq_{\text{temp}} = \freq_{\text{La}} \cdot 2^{n/12}
\label{eq:temperamento_igual}
\end{equation}

donde:
\begin{itemize}
    \item $\freq_{\text{La}}$ es la frecuencia del La de referencia (por defecto 415 Hz para barroco)
    \item $n$ es el número de semitonos desde La
\end{itemize}

Para cada nota:
\begin{align}
\text{Re} \quad (n = -7): \quad &\freq_{\text{temp}} = 415 \cdot 2^{-7/12} \approx \SI{311.13}{\hertz} \\
\text{Mi} \quad (n = -5): \quad &\freq_{\text{temp}} = 415 \cdot 2^{-5/12} \approx \SI{349.23}{\hertz} \\
\text{Fa\#} \quad (n = -3): \quad &\freq_{\text{temp}} = 415 \cdot 2^{-3/12} \approx \SI{392.00}{\hertz} \\
\text{La} \quad (n = 0): \quad &\freq_{\text{temp}} = 415 \cdot 2^{0/12} = \SI{415.00}{\hertz}
\label{eq:frecuencias_notas}
\end{align}

\subsubsection{Cálculo de Desviación}

La desviación en cents se calcula como:

\begin{align}
\text{Desviación} &= 1200 \log_2\left(\frac{\freqzero}{\freq_{\text{temp}}}\right) \\
&= 1200 \left[\log_2(\freqzero) - \log_2(\freq_{\text{temp}})\right] \\
&= 1200 \left[\log_2(\freqzero) - \log_2(\freq_{\text{La}}) - \frac{n}{12}\right]
\label{eq:desviacion_deriv}
\end{align}

\subsubsection{Interpretación}

\begin{itemize}
    \item \textbf{Desviación = 0 cents}: Frecuencia perfectamente temperada
    \item \textbf{Desviación > 0 cents}: Nota más aguda que el temperamento
    \item \textbf{Desviación < 0 cents}: Nota más grave que el temperamento
\end{itemize}

Esta métrica es fundamental para evaluar la afinación del instrumento respecto a estándares musicales.

\subsection{Altura de Picos de Admitancia}

\subsubsection{Definición}

La altura del pico de admitancia en la frecuencia de resonancia $\freqzero$ es:

\begin{equation}
\admittance_{\max} = \max_f |\admittance(f)|
\label{eq:admittance_peak}
\end{equation}

donde la admitancia es:
\begin{equation}
\admittance(f) = \frac{1}{\impedance(f)}
\label{eq:admittance_def}
\end{equation}

\subsubsection{Relación con Facilidad de Emisión}

La altura del pico está relacionada con la facilidad de emisión del sonido. Un pico más alto indica:

\begin{itemize}
    \item Mayor transferencia de energía desde el músico al instrumento
    \item Menor impedancia de entrada
    \item Mayor eficiencia acústica
\end{itemize}

Físicamente, esto se relaciona con:
\begin{equation}
P_{\text{acústica}} = \frac{1}{2} |U|^2 \Re(\impedance) = \frac{1}{2} |p|^2 \Re(\admittance)
\label{eq:potencia_acustica}
\end{equation}

donde $P_{\text{acústica}}$ es la potencia acústica generada.

\subsection{Ratio de Armónicos Pares/Impares}

\subsubsection{Definición}

El ratio de armónicos pares/impares caracteriza el contenido espectral del sonido:

\begin{equation}
R = \frac{A_{\text{pares}}}{A_{\text{impares}}} = \frac{A_2 + A_4 + A_6 + \ldots}{A_1 + A_3 + A_5 + \ldots}
\label{eq:harmonic_ratio}
\end{equation}

donde $A_n$ es la amplitud del $n$-ésimo armónico.

\subsubsection{Cálculo desde Admitancia}

Las amplitudes se extraen de los picos de admitancia en las frecuencias de antiresonancia:

\begin{align}
A_1 &= |\admittance(\freqzero)| \\
A_2 &= |\admittance(\freqone)| \\
A_3 &= |\admittance(\freqtwo)|
\label{eq:amplitudes_armonicos}
\end{align}

donde $\freqtwo$ es la tercera frecuencia de antiresonancia.

El ratio se calcula como:
\begin{equation}
R = \frac{A_2}{A_1 + A_3}
\label{eq:ratio_calculo}
\end{equation}

\subsubsection{Interpretación Musical}

\begin{itemize}
    \item \textbf{$R > 1$}: Predominio de armónicos pares, sonido más brillante
    \item \textbf{$R < 1$}: Predominio de armónicos impares, sonido más oscuro
    \item \textbf{$R \approx 1$}: Balance entre armónicos, sonido equilibrado
\end{itemize}

\subsection{Coherencia de Fase}

\subsubsection{Definición}

La coherencia de fase mide la diferencia de fase entre armónicos:

\begin{equation}
\Delta\phi = \arg(\impedance(\freqone)) - \arg(\impedance(\freqzero))
\label{eq:phase_coherence}
\end{equation}

donde $\arg(\impedance)$ es la fase de la impedancia compleja.

\subsubsection{Relación con Claridad del Sonido}

La fase de la impedancia está relacionada con la respuesta temporal del instrumento. Una mayor coherencia de fase indica:

\begin{itemize}
    \item Mayor sincronización entre armónicos
    \item Sonido más claro y definido
    \item Menor distorsión temporal
\end{itemize}

Físicamente, la fase está relacionada con el retardo de grupo:
\begin{equation}
\tau_g = -\frac{d\phi}{d\omega}
\label{eq:retardo_grupo}
\end{equation}

\subsection{Estabilidad de Pitch}

\subsubsection{Definición}

La estabilidad de pitch se basa en la pendiente de la fase cerca de la resonancia:

\begin{equation}
S = \left|\frac{d\phi}{df}\right|_{f=\freqzero}
\label{eq:pitch_stability}
\end{equation}

\subsubsection{Derivación}

Para un resonador, la fase de la impedancia cerca de la resonancia puede aproximarse como:

\begin{equation}
\phi(f) \approx \arctan\left(Q\left(\frac{f}{\freqzero} - \frac{\freqzero}{f}\right)\right)
\label{eq:fase_aproximacion}
\end{equation}

La derivada es:
\begin{equation}
\frac{d\phi}{df} = \frac{Q/\freqzero}{1 + Q^2\left(\frac{f}{\freqzero} - \frac{\freqzero}{f}\right)^2}
\label{eq:derivada_fase}
\end{equation}

Evaluando en $f = \freqzero$:
\begin{equation}
S = \left|\frac{d\phi}{df}\right|_{f=\freqzero} = \frac{Q}{\freqzero}
\label{eq:estabilidad_final}
\end{equation}

\subsubsection{Interpretación}

\begin{itemize}
    \item \textbf{$S$ alto}: Mayor estabilidad, menor sensibilidad a variaciones de presión
    \item \textbf{$S$ bajo}: Menor estabilidad, mayor sensibilidad a variaciones
\end{itemize}

\subsection{Frecuencia de Corte (Cut-off Frequency)}

\subsubsection{Definición}

La frecuencia de corte es la frecuencia más alta donde la admitancia normalizada cae por debajo de un umbral:

\begin{equation}
\freq_{\text{cut}} = \max\left\{f : \frac{|\admittance(f)|}{|\admittance_{\max}|} \geq \theta\right\}
\label{eq:cutoff_freq}
\end{equation}

donde $\theta$ es el umbral (típicamente 0.1 = 10\%).

\subsubsection{Interpretación Física}

La frecuencia de corte está relacionada con:

\begin{itemize}
    \item Límite superior del rango útil del instrumento
    \item Capacidad de generar armónicos superiores
    \item Eficiencia de radiación a altas frecuencias
\end{itemize}

Físicamente, está relacionada con el diámetro del conducto:
\begin{equation}
\freq_{\text{cut}} \propto \frac{c_0}{d}
\label{eq:cutoff_diametro}
\end{equation}

donde $d$ es el diámetro característico del conducto.

% ==================== SECCIÓN 5: VISUALIZACIONES ====================
\section{Visualizaciones}

\subsection{Visualización 2D}

El sistema proporciona múltiples tipos de visualización 2D:

\subsubsection{Perfil Físico}

Muestra el ensamblaje real de la flauta con uniones mortise/tenon, respetando las longitudes físicas de cada parte.

\begin{figure}[H]
\centering
\includegraphics[width=0.9\columnwidth]{figuras/perfil_fisico.png}
\caption{Perfil físico mostrando ensamblaje real}
\label{fig:perfil_fisico}
\end{figure}

\subsubsection{Perfil Acústico}

Muestra la geometría acústica concatenada, con el corcho (stopper) en $x=0$, mostrando la longitud acústica efectiva.

\begin{figure}[H]
\centering
\includegraphics[width=0.9\columnwidth]{figuras/perfil_acustico.png}
\caption{Perfil acústico con corcho en x=0}
\label{fig:perfil_acustico}
\end{figure}

\subsubsection{Vista Sólido 2D}

Muestra el perfil interno y externo con el espesor de pared, permitiendo visualizar la geometría completa del instrumento.

\begin{figure}[H]
\centering
\includegraphics[width=0.9\columnwidth]{figuras/vista_solido_2d.png}
\caption{Vista sólido 2D mostrando espesor de pared}
\label{fig:vista_solido_2d}
\end{figure}

\subsubsection{Corte Axial}

Muestra un corte del sólido en el plano axial, visualizando los agujeros como cilindros o conos según su configuración.

\begin{figure}[H]
\centering
\includegraphics[width=0.9\columnwidth]{figuras/corte_axial.png}
\caption{Corte axial del sólido mostrando agujeros}
\label{fig:corte_axial}
\end{figure}

\subsection{Visualización 3D}

El sistema genera modelos 3D paramétricos usando CadQuery y los visualiza con PyVista.

\begin{figure}[H]
\centering
\includegraphics[width=0.9\columnwidth]{figuras/flauta_3d.png}
\caption{Modelo 3D de flauta completa}
\label{fig:flauta_3d}
\end{figure}

\subsection{Gráficos de Análisis Acústico}

\subsubsection{Admitancia en Frecuencia}

Muestra la admitancia acústica en función de la frecuencia, con valores de los primeros picos anotados.

\begin{figure}[H]
\centering
\includegraphics[width=0.9\columnwidth]{figuras/admitancia_frecuencia.png}
\caption{Admitancia en función de la frecuencia}
\label{fig:admitancia_frecuencia}
\end{figure}

\subsubsection{Presión y Flujo Espaciales}

Muestra la distribución espacial de presión y flujo acústicos, sincronizados con la geometría del instrumento.

\begin{figure}[H]
\centering
\includegraphics[width=0.9\columnwidth]{figuras/presion_flujo_espacial.png}
\caption{Distribución espacial de presión y flujo}
\label{fig:presion_flujo}
\end{figure}

% ==================== SECCIÓN 6: CASOS DE ESTUDIO ====================
\section{Casos de Estudio}

\subsection{Caso 1: Análisis de Flauta Histórica}

[Incluir análisis detallado de una flauta específica con capturas de pantalla]

\subsection{Caso 2: Comparación de Dos Flautas}

[Incluir comparación lado a lado con análisis de diferencias]

\subsection{Caso 3: Optimización de Diseño}

[Incluir ejemplo de modificación de geometría y evaluación de impacto acústico]

% ==================== SECCIÓN 7: BASE DE DATOS ====================
\section{Base de Datos}

\subsection{Esquema}

La base de datos SQLite almacena:

\begin{itemize}
    \item \textbf{Tabla \texttt{flute}}: Información general de cada flauta
    \item \textbf{Tabla \texttt{flute\_geometry}}: Mediciones geométricas discretas
    \item \textbf{Tabla \texttt{calculation\_parameters}}: Parámetros de cálculo (temperatura, diapasón, etc.)
    \item \textbf{Tabla \texttt{acoustic\_analysis\_results}}: Resultados de impedancia serializados
    \item \textbf{Tabla \texttt{external\_geometry}}: Perfiles externos
\end{itemize}

\subsection{Optimización}

Para reducir el tamaño de la base de datos:

\begin{itemize}
    \item Caché de resultados de impedancia (evita recálculos)
    \item Hash de parámetros para detección de cambios
    \item Optimización de datos de presión/flujo (solo primeros 3 armónicos)
    \item Limpieza periódica de datos grandes
\end{itemize}

% ==================== SECCIÓN 8: HERRAMIENTAS ADICIONALES ====================
\section{Herramientas Adicionales}

\subsection{Editor de Geometría}

Editor interactivo que permite:
\begin{itemize}
    \item Modificar puntos del perfil interno (arrastrar, añadir, eliminar)
    \item Editar agujeros (posición, diámetro, altura de chimenea)
    \item Control de versiones (undo/redo)
    \item Comparación acústica original vs modificado
    \item Guardado como nueva flauta
\end{itemize}

\subsection{Generación de G-code}

Genera código NC para torno CNC:
\begin{itemize}
    \item Estrategias de desbaste (longitudinal/transversal)
    \item Parámetros configurables (velocidad, avance, profundidad)
    \item Visualización de trayectorias
    \item Modo económico opcional
\end{itemize}

\subsection{Planos de Ingeniería}

Genera PDFs vectoriales profesionales:
\begin{itemize}
    \item Múltiples páginas (partes individuales + ensamblado)
    \item Tablas de dimensiones y agujeros
    \item Grilla milimétrica
    \item Cotas y anotaciones
\end{itemize}

% ==================== REFERENCIAS ====================
\begin{thebibliography}{99}

\bibitem{openwind2020}
Helie, T., \& Matignon, D. (2020).
\textit{OpenWind: A Python Library for Wind Instrument Modeling}.
INRIA Research Report.

\bibitem{webster1919}
Webster, A. G. (1919).
Acoustical impedance, and the theory of horns and of the phonograph.
\textit{Proceedings of the National Academy of Sciences}, 5(7), 275-282.

\bibitem{benade1990}
Benade, A. H. (1990).
\textit{Fundamentals of Musical Acoustics}.
Dover Publications, New York.

\bibitem{chaigne2013}
Chaigne, A., \& Kergomard, J. (2013).
\textit{Acoustics of Musical Instruments}.
Springer, New York.

\bibitem{flutehistory}
Toff, N. (1996).
\textit{The Flute Book: A Complete Guide for Students and Performers}.
Oxford University Press, Oxford.

\bibitem{quantz1752}
Quantz, J. J. (1752).
\textit{Versuch einer Anweisung die Flöte traversiere zu spielen}.
Berlin.

\bibitem{ricci2003}
Ricci, A. (2003).
\textit{The Baroque Flute Fingering Book}.
Ricci Publications.

\bibitem{powell2002}
Powell, A. (2002).
\textit{The Flute}.
Yale University Press, New Haven.

\end{thebibliography}

\end{document}

